\lesson{5}{Oct 11 2021 Mon (11:28:43)}{Federalism}{Unit 2}

\subsubsection*{The types of Power in the Constitution}

The \textbf{U.S. Constitution} lists the delegated powers. They are the powers that are assigned to the \textbf{National Government}. They are also known as the \textit{Expressed Powers} because they are directly named in the Constitution. The Constitution and Amendments also include some of the reserved powers. They are not given to the National Government. They are used by the State's Government. The National and State Government share powers called the \textbf{Concurrent Powers}.

The \textbf{Constitution} also has implied powers. American leaders like Thomas Jefferson argued that the \textbf{"Elastic Clause"} in the Constitution gave too much power to Congress by giving the ability to make all of the laws. They feared that the clause implied Congress had powers beyond those directly named.

The issue of using Implied Powers was argued during the Supreme Court case of \textbf{McCulloch v. Maryland}. The U.S. Government established a national bank in \textbf{Maryland}. Then, the state passed legislation imposing taxes on the bank. Then, the cashier of the Baltimore Branch of the bank, \textit{James W. McCulloch} refused to pay the taxes.

This brought up questions for the government:

\begin{question}
    Does the Constitution allow the national government to run a bank?
\end{question}

\begin{answer}
    Yes. It doesn't say so directly in the Constitution. But, still, the court ruled it an implied power. Congress had hoped that they could set up a single currency and a national bank to help resolve some interstate commerce challenges. Even though Congress never set up its own business corporation before, they deemed it \textbf{"Necessary and Proper} to be carrying out the duties of regulating interstate commerce. Today, the national government runs several corporations, such as AMTRAK, the passenger railroad system.
\end{answer}

\begin{question}
    Does the Constitution allow state governments to tax a national bank operating within its borders?
\end{question}

\begin{answer}
    No. \textit{Justice Marshall} was very clear that this could undermine every aspect of the national government. If a state could tax the national bank, then that could be a precedent for it to tax the mail service. They could probably even tax it so much that there would be no funds so the mail can operate.
\end{answer}

Let's quickly go over the powers again:

\begin{definition}[Delegated to the National Government]
    The power to:
    
    \begin{enumerate}
        \item Organize and control trade between states and with other nations
        \item Coin and print money
        \item Make agreements with foreign nations
        \item Establish post offices and interstate highways
        \item Raise and support military forces
        \item Declare war and make peace
        \item Govern U.S. territories and admit new states, and create rules and laws for people who wish to move into the nation and about how they may become citizens
    \end{enumerate}
\end{definition}

\begin{definition}[Concurrent or Shared Powers]
    The power to:
        
    \begin{enumerate}
        \item Raise and collect taxes
        \item Borrow money
        \item Establish courts
        \item Grant permissions to start banks
        \item Enforce laws and punish lawbreakers
        \item Provide health care services and other assistance to the people
    \end{enumerate}
\end{definition}    

\begin{definition}[Reserved to the State Governments]
    The power to:
    
    \begin{enumerate}
        \item Create local government
        \item Set up businesses
        \item Build and support public schools
        \item Organize and control trade within the state
        \item Conduct elections
        \item Determine rules for voting, such as living in the district where the person votes
        \item Make marriage laws
        \item Determine the requirements for professional workers, like doctors and teachers
    \end{enumerate}
\end{definition}

\subsubsection*{Constitution Addressing Federalism}

\newpage
